\subsection{The Merritt's Implementation}

\subsubsection{Implementation}

This approach utilizes the pandas Python library for data processing, reading in data as a tabular format and applying operations to entire columns at a time to achieve faster transformations. The implementation and can be divided into four parts: initial data processing, model generation, testing against a model, and jackknife training. Each part is implemented in a separate function to allow for flexibility in terms of $n$-gram size, excluding initial data processing which only occurs once and does not take any arguments. 

In in the initial data processing stage, the tagged sentence data is read from the spreadsheet, with each row containing a single word in a sentence and the expected part of speech, and with numbers marking the beginning of sentences. To simplify the process of dividing the dataset into sentences for training and testing, the sentence numbers are cascaded down to the subsequent words that occur before the next sentence starts. Additionally, the sentence number is extracted from the original format by stripping away non-numeric characters. Finally, string data types needed to be specified and missing value detection needed to be disabled when loading the dataset in using pandas to avoid problems with empty fields or words reserved for values in Python, i.e. None and False.

In the model generation stage, two major steps occur: mapping and flattening. Before mapping can occur, the history of previous parts of speech must be generated for each word and stored in the row for later use. To do this, the part of speech column is copied and shifted down based on the number of history entries required. For example, when generating quadgrams, the history requires the previous three parts of speech, meaning that three shifted columns need to be added. Values in the shifted columns are left blank if the sentence number changes after a shift. With the history added to each row in the table, a mapping can be generated by applying a function to the rows of the table, creating a dictionary that maps history patterns to occurrence counters for each part of speech. This step of applying the function to the rows to build the dictionary is the most time consuming step in our implementation, taking roughly 15 seconds per jackknife iteration. With the mapping dictionary created, a flattening step is performed to create the model where the mappings are separated into separate dictionaries for unanimous and non-unanimous n-grams that skip past the counts for each part of speech and instead return only the most probable part of speech. 

One modification explored but ultimately not implemented at this model generation stage was aimed at reducing ambiguity, but it was not included due to negative impact on accuracy. Specifically, in the case of the non-unanimous model, we explored removing cases where the part of speech decision was ambiguous. For example, when using the first thousand sentences as the testing set and the remaining sentences as the training set, one such ambiguity occurs for the the word 'ruling' after a preposition, proper noun, then possesive ending. In this case, there is an equal chance of 'ruling' representing a 'Verb, gerund or present participle' ('VBG') or a 'Noun, singular or mass' ('NN'), with both occurring 5 times in the training set. Removing this pattern from our model did not cause a noticeable increase in accuracy, so ambiguity is instead currently decided on a seemingly arbitrary basis. Currently, in the case of ties, the pattern that was observed first will be selected by the max function call on the dictionary of counts.

\begin{lstlisting}[language=Python]
    def generate_ngram_model(training_set, gram_size=4):
        # ---------------- STAGE 0: PREPROCESSING HISTORY ----------------
        modified = pd.DataFrame(training_set)
    
        for n in range(gram_size - 1):
            modified[f'Previous POS {n + 1}'] = modified['POS'].shift(n + 1)
            modified[f'Previous POS {n + 1}'] = modified[f'Previous POS {n + 1}'].where(modified['Sentence #'] == modified['Sentence #'].shift(n + 1))
    
        # ---------------- STAGE 1: INITIAL COUNTING ----------------
        ngram_map = {}
    
        def write_to_map(row):
            previous_pos = list(map(lambda n : row[f'Previous POS {n}'], list(range(1, gram_size))[::-1]))
            previous_pos = list(filter(lambda pos : not pd.isna(pos), previous_pos))
    
            patterns = [(*previous_pos[-i:], row['Word']) for i in range(1, len(previous_pos) + 1)]
            patterns.append(row['Word'])
    
            for pattern in patterns:
                if pattern in ngram_map:
                    if row['POS'] in ngram_map[pattern]:
                        ngram_map[pattern][row['POS']] += 1
                    else:
                        ngram_map[pattern][row['POS']] = 1
                else:
                    ngram_map[pattern] = { row['POS']: 1 }
    
        modified.apply(write_to_map, axis=1)
    
        # ---------------- STAGE 2: FLATTENING ----------------
    
        # index 0 is the unanimous map, index 1 is the highest probability map
        ngram_model = [{}, {}]
        for (key, pos_map) in ngram_map.items():
            map_index = 0 if len(pos_map) == 1 else 1
    
            ngram_model[map_index][key] = max(pos_map, key=pos_map.get)
    
        return (ngram_map, ngram_model)
\end{lstlisting}

In the testing against a model stage, the model from the previous stage is taken alongside the testing set, and each row is tested using both the unanimous and non-unanimous models when necessary. First the testing set is augmented with history similar to the training set, and then a function is applied on each row that tests and stores the outcomes into a new column. The testing order is to go through all the unanimous models from least amount of context (unigram) to most amount of context (quadgram in our case), then go through all the non-unanimous models in the opposite order. Output is appended to a text file for debugging and logging. This log file should be deleted between separate iterations of the program to prevent confusion.

In the final jackknifing stage, the raw dataset is divided into training and testing sets that can be passed to the model generation and testing stages. Each iteration selects 1000 sentences for testing and passes the remaining sentences as the training set. Between iterations, the entire model is discarded and rebuilt, with only the correct and incorrect result counts remaining across iterations. Overall accuracy is output at the end.

\subsubsection{Results}
Ultimately, this approach yielded lower than expected accuracy and generally performed slower than other approaches. In terms of the performance, the biggest impact was in the model generation stage when the mapping function is applied to the rows of the training set. We were able to iterate on this, having initially iterated over the rows individually, taking 2-3 minutes per training set. In the end, we reduced this down to about 15 seconds by applying a function to the rows of the table instead, but this still adds up. Regarding the lower accuracy, comparing our quadgram results to the provided sample output (99.00\% accuracy), we found that only additional incorrect tests we had were when ambiguous decisions occurred, making up the 4\% difference in accuracy. 

The results of this implementation of the $n$-gram tagger were as follows:
\begin{itemize}
    \item \textbf{Unigram only:} 93.60\% accuracy, runtime of 5m 13.980s
    \item \textbf{Up to bigram:} 95.84\% accuracy, runtime of 7m 46.116s
    \item \textbf{Up to trigram:} 95.94\% accuracy, runtime of 10m 46.101s
    \item \textbf{Up to quadgram:} 95.91\% accuracy, runtime of 13m 39.980s
\end{itemize}